% ============================================================
%  Dispensa 03: Curva EIOPA Risk-Free — Metodo Smith-Wilson
%               Approccio variazionale (struttura stile PCA)
%  Corso di Calcolo Numerico — Laboratorio
%  Università degli Studi di Verona — A.A. 2026–2027
% ============================================================
\documentclass[11pt,a4paper]{article}

\usepackage[T1]{fontenc}
\usepackage[utf8]{inputenc}
\usepackage[italian]{babel}
\usepackage{lmodern}
\usepackage{amsmath,amssymb,amsthm,mathtools}
\usepackage{geometry}
\usepackage{booktabs}
\usepackage{array}
\usepackage{enumitem}
\usepackage{graphicx}
\usepackage[table]{xcolor}
\usepackage{hyperref}
\usepackage{float}
\usepackage{fancyvrb}
\usepackage{tcolorbox}
\tcbuselibrary{breakable}
\usepackage{titlesec}

% le figure sono prodotte da R/eiopa_rfr_sw_variazionale.R in output/eiopa_rfr_sw_var/
\graphicspath{{}}

\geometry{top=2.5cm,bottom=2.5cm,left=2.8cm,right=2.8cm}
\setlist[itemize]{topsep=4pt,itemsep=2pt,parsep=0pt}
\setlist[enumerate]{topsep=4pt,itemsep=2pt,parsep=0pt}
\hypersetup{colorlinks=true,linkcolor=blue!55!black,urlcolor=blue!55!black}

\titleformat{\section}{\large\bfseries}{\thesection.}{0.55em}{}[{\vspace{1pt}\titlerule[0.4pt]}]
\titleformat{\subsection}{\normalsize\bfseries}{\thesubsection.}{0.55em}{}
\titleformat{\subsubsection}{\small\bfseries\itshape}{\thesubsubsection.}{0.5em}{}

%--- ambienti teorema (solo enunciati, niente dimostrazioni)
\theoremstyle{plain}
\newtheorem{teorema}{Teorema}[section]
\newtheorem{proposizione}[teorema]{Proposizione}
\theoremstyle{definition}
\newtheorem{definizione}[teorema]{Definizione}
\newtheorem{esempio}[teorema]{Esempio}
\theoremstyle{remark}
\newtheorem*{osservazione}{Osservazione}
\newtheorem*{attenzione}{Nota numerica}

%--- macro acronimi (funzionano dentro e fuori math mode)
\newcommand{\UFR}{\ensuremath{\mathrm{UFR}}}
\newcommand{\LLP}{\ensuremath{\mathrm{LLP}}}
\newcommand{\CP}{\ensuremath{\mathrm{CP}}}
\newcommand{\CRA}{\ensuremath{\mathrm{CRA}}}
\newcommand{\Qmat}{\mathbf{Q}}
\newcommand{\Hmat}{\mathbf{H}}
\newcommand{\transp}{^{\!\top}}

%--- box colorati
\tcbset{
  mybase/.style={
    colback=gray!5, colframe=gray!40,
    boxrule=0.4pt, arc=2pt,
    left=7pt, right=7pt, top=5pt, bottom=5pt,
    fonttitle=\small\bfseries
  }
}
\newtcolorbox{keybox}[1][Concetto chiave]{mybase,
  colframe=blue!40!black, colback=blue!4, title={#1}}
\newtcolorbox{algobox}[1][Algoritmo]{mybase,
  colframe=teal!50!black, colback=teal!4, title={#1}}
\newtcolorbox{warningbox}{mybase,
  colframe=orange!70!black, colback=orange!5}
\newtcolorbox{codebox}[1][Codice R]{mybase,
  colframe=green!40!black, colback=green!3, title={#1}}

% ============================================================
\title{%
  \textbf{Costruzione della Curva EIOPA Risk-Free}\\[6pt]
  \large Il metodo Smith--Wilson: un approccio variazionale\\[4pt]
}
\author{Laboratorio di Calcolo Numerico --- Laurea Triennale in Matematica Applicata\\
  Universit\`a degli Studi di Verona\\[2pt]
  Anno Accademico 2026--2027}
\date{}

\begin{document}
\maketitle
\tableofcontents
\newpage

% ============================================================
\section{Motivazioni economiche}
% ============================================================

\subsection{Il problema: valutare oggi impegni decennali}

Una compagnia di assicurazione vita o un fondo pensione assume \textbf{obblighi di pagamento
che si estendono per decenni}: una rendita vitalizia o una polizza possono generare uscite
fino a 50--60 anni nel futuro.

Il quadro regolamentare europeo \textbf{Solvency~II} chiede che ogni compagnia dimostri,
\emph{oggi}, di avere risorse sufficienti per far fronte a tutti questi impegni futuri ---
\`e il principio di \textbf{solvibilit\`a}. Per farlo occorre rispondere alla domanda:
\emph{quanto valgono oggi quei pagamenti futuri?}

La risposta dipende dal \textbf{valore temporale del denaro}: 100\,\texteuro{} ricevuti tra
30 anni valgono \emph{meno} di 100\,\texteuro{} ricevuti oggi, perch\'e nel frattempo il
denaro pu\`o essere investito e generare interessi. Quanto meno? Dipende dal
\textbf{tasso di interesse} applicabile su quell'orizzonte temporale.

Poich\'e le passivit\`a hanno scadenze diverse, serve conoscere un tasso di interesse per
ogni orizzonte: 1 anno, 5 anni, 20 anni, 50 anni\ldots\ L'insieme di questi tassi al
variare della scadenza $T$ \`e la \textbf{struttura a termine dei tassi di interesse}, o
pi\`u brevemente \textbf{curva dei tassi}.

Solvency~II impone che tutte le imprese europee usino la \textbf{stessa} curva, calcolata e
pubblicata mensilmente da \textbf{EIOPA} (European Insurance and Occupational Pensions
Authority). La metodologia ufficiale \`e descritta nel documento tecnico
EIOPA-BoS-25-599~\cite{eiopa2025}.

\subsection{Il problema dei dati: il mercato finisce a 20 anni}

Il mercato finanziario fornisce quotazioni di tassi di interesse per molte scadenze tramite
contratti standardizzati. Per l'euro questi contratti sono \textbf{attivamente scambiati e
affidabili} fino a circa 20 anni di scadenza: questo orizzonte \`e detto
\textbf{Last Liquid Point} ($\LLP = 20$ anni). Oltre tale soglia le quotazioni diventano
sporadiche, poco rappresentative o del tutto assenti.

Il problema \`e evidente: le passivit\`a assicurative si estendono ben oltre i 20 anni,
ma i dati di mercato si esauriscono proprio l\`i. EIOPA deve comunque fornire una stima
del tasso di interesse anche per le scadenze lunghe, fino al \textbf{Convergence Point}
($\CP = 60$ anni).

\begin{keybox}[Due problemi in uno]
Costruire la curva EIOPA significa risolvere due problemi distinti:
\begin{enumerate}
  \item \textbf{Ricostruire} in modo regolare i tassi alle scadenze dove esistono
        quotazioni di mercato ($T \le \LLP = 20$ anni);
  \item \textbf{Estendere} la curva oltre il $\LLP$ fino al $\CP = 60$ anni in modo
        stabile, coerente e non arbitrario.
\end{enumerate}
\end{keybox}

\subsection{L'ancora di lungo periodo: l'UFR}

Come estendere la curva oltre i 20 anni? L'idea ingenua --- prolungare l'ultimo tasso
osservato ``a piatto'' --- \`e fragile: una piccola variazione nel tasso a 20 anni si
propagherebbe su tutti gli anni da 21 a 60, amplificando le oscillazioni nei bilanci
assicurativi.

La scelta regolatoria (\cite{eiopa2025}, sez.~3.6--3.8) \`e di \textbf{agganciare la curva
a un tasso di equilibrio di lungo periodo}: l'\textbf{Ultimate Forward Rate}
($\UFR = 3{,}30\%$ per l'euro). L'$\UFR$ rappresenta la stima regolamentare del tasso
reale di crescita economica di lungo periodo sommato all'inflazione attesa; cambia molto
lentamente nel tempo ed \`e quindi un'\textbf{ancora stabile}.

La curva viene costruita in modo che il tasso di interesse, superato il $\LLP$, converga
\emph{gradualmente e in modo regolare} verso l'$\UFR$ entro il $\CP = 60$ anni.

\begin{figure}[H]
\centering
\includegraphics[width=0.92\textwidth]{../output/eiopa_rfr_sw_var/fig1_motivazione_estrapolazione.pdf}
\caption{La curva EUR calcolata da EIOPA (dicembre~2025) fino al Convergence Point di 60 anni.
Nella zona ombreggiata ($T > \LLP = 20$ anni) non esistono dati di mercato affidabili: la
curva \`e costruita in modo che il tasso converga al valore di lungo periodo
$\UFR = 3{,}30\%$. \`E questo l'obiettivo del metodo Smith--Wilson.}
\label{fig:motivazione}
\end{figure}

\subsection{Eliminare il rischio di controparte: il CRA}
\label{sec:cra}

I contratti di mercato usati come sorgente dati (swap di tasso di interesse) includono un
piccolo \textbf{premio per il rischio di controparte}: chi riceve i pagamenti sopporta il
rischio che la controparte non onori il contratto. Questo premio non riflette il valore
temporale del denaro in s\'e, e deve essere rimosso.

Per ottenere un tasso \emph{privo di rischio} EIOPA sottrae un \textbf{Credit Risk
Adjustment} ($\CRA = 10$ bps per l'euro, \cite{eiopa2025} sez.~3.5) dai tassi di mercato
\emph{prima} di calibrare la curva. Tutti i tassi di input che useremo nella
Sez.~\ref{sec:ricostruzione} sono gi\`a corretti di questi 10 punti base.

% ============================================================
\section{La curva dei tassi e i tassi forward}
% ============================================================

Tutte le nozioni di tasso --- spot, forward, par --- sono contenute in un unico oggetto:
il \textbf{fattore di sconto} $P(0,t)$. Spot e forward sono semplicemente due modi diversi
di ``leggere'' questa stessa funzione. Costruiamo la sua definizione partendo dalla
capitalizzazione discreta e arrivando alla capitalizzazione continua come caso limite naturale.

% ------------------------------------------------------------
\subsection{Dalla capitalizzazione al fattore di sconto}
% ------------------------------------------------------------

\paragraph{Capitalizzazione annua composta.}
\emph{Capitalizzare} significa trasformare un euro di oggi in pi\`u euro in futuro.
Investendo $1$~\texteuro{} a tasso $r$ per $T$ anni si ottiene il
\textbf{fattore di capitalizzazione}
\[
  C_{\rm ann}(r,T) = (1+r)^T.
\]
Ad esempio con $r=3\%$ e $T=10$ anni: $C = 1.03^{10} \approx 1.344$, cio\`e $1$~\texteuro{}
cresce a $1.34$~\texteuro{}.

\paragraph{Capitalizzazione $m$ volte l'anno.}
Se la capitalizzazione avviene $m$ volte per anno al tasso nominale $r_m$ (rate $r_m/m$
per periodo), il fattore di capitalizzazione diventa
\[
  C_m(r_m,T) = \left(1 + \frac{r_m}{m}\right)^{mT}.
\]
All'aumentare di $m$ la capitalizzazione diventa sempre pi\`u frequente: $m=12$ mensile,
$m=365$ giornaliera, e cos\`i via.

\paragraph{Limite continuo: $m \to \infty$.}
Usando il limite fondamentale $\lim_{m\to\infty}(1+x/m)^m = e^x$ si ottiene il
\textbf{fattore di capitalizzazione continua}:
\begin{equation}
  C_\infty(r,T) = \lim_{m\to\infty}\left(1+\frac{r_m}{m}\right)^{mT} = e^{r\,T}.
  \label{eq:cap-cont}
\end{equation}
La capitalizzazione continua \`e lo standard della matematica finanziaria quantitativa
perch\'e l'esponenziale rende i \emph{logaritmi} dei fattori additivi (invece che i fattori
stessi moltiplicativi), aprendo la strada alla derivata che definisce il tasso forward.

\paragraph{Il fattore di sconto: il problema inverso.}
Finora abbiamo chiesto: \emph{``quanto vale in futuro un euro investito oggi?''} Il problema
\emph{inverso} --- \emph{``quanto vale oggi un euro che riceveremo al tempo $t$?''} --- \`e
il problema dello sconto. Il fattore di sconto \`e semplicemente il reciproco del fattore
di capitalizzazione continua.

\begin{definizione}[Fattore di sconto]
\label{def:sconto}
Il \textbf{fattore di sconto} $P(s,t)$ \`e il valore al tempo $s$ di $1$~\texteuro{}
pagato al tempo $t \ge s$:
\[
  P(s,t) = \frac{1}{C_\infty(r,\,t-s)} = e^{-r(s,t)\,(t-s)}.
\]
Il primo argomento $s$ \`e la \textbf{data di osservazione} (quando valutiamo), il secondo
$t$ \`e la \textbf{data di pagamento} (quando riceviamo l'euro). Nel seguito fissiamo
$s = 0$ (oggi) e scriviamo semplicemente $P(0,t)$.

Propriet\`a immediate:
\begin{itemize}
  \item $P(0,0) = 1$ (un euro pagato subito vale un euro);
  \item $P(0,t)$ strettamente decrescente in $t$ se $r > 0$;
  \item tutta la struttura a termine dei tassi \`e codificata in $P(0,\cdot)$.
\end{itemize}
\end{definizione}

% ------------------------------------------------------------
\subsection{Tasso spot: leggere $P(0,t)$}
% ------------------------------------------------------------

Dalla Definizione~\ref{def:sconto} il fattore di sconto ha la forma $P(0,t)=e^{-r(t)t}$.
Invertendo si ricava il tasso che lo genera.

\begin{definizione}[Tasso spot]
\label{def:spot}
Il \textbf{tasso spot} $r(t)$ \`e il tasso in capitalizzazione continua associato al
fattore di sconto $P(0,t)$:
\begin{equation}
  r(t) = -\frac{\ln P(0,t)}{t}.
  \label{eq:spot-inv}
\end{equation}
\`E la media logaritmica del costo del denaro sull'intero intervallo $[0,t]$.
\end{definizione}

\paragraph{Perch\'e il logaritmo rende tutto pi\`u semplice.}
In capitalizzazione discreta i fattori di sconto si \emph{moltiplicano}: il fattore su
$[0,t_1+t_2]$ \`e il prodotto di quello su $[0,t_1]$ e di quello su $[t_1,t_1+t_2]$,
il che rende le composizioni scomode da derivare. In capitalizzazione continua, invece,
$\ln P(0,t) = -r(t)\cdot t$ \`e \emph{additivo}: i log-fattori si sommano, e questo
\`e il punto di ingresso naturale per la derivata che definisce il forward.

\paragraph{Conversione con la capitalizzazione annua.}
Il tasso annuo composto $r_{\rm ann}(t) = P(0,t)^{-1/t}-1$ \`e legato al tasso continuo da
\[
  r_{\rm ann}(t) = e^{r(t)} - 1, \qquad r(t) = \ln(1 + r_{\rm ann}(t)).
\]
EIOPA pubblica i risultati in capitalizzazione annua (i par swap pagano cedole annue),
ma il metodo Smith--Wilson costruisce $P(0,t)$ internamente usando la capitalizzazione
continua: la conversione $r_{\rm ann} = e^r - 1$ avviene solo all'output finale.

% ------------------------------------------------------------
\subsection{Tasso forward: dall'intervallo finito al limite istantaneo}
% ------------------------------------------------------------

\paragraph{Forward su intervallo finito: assenza di arbitraggio.}
Il \textbf{tasso forward} $f(t_1, t_2)$ \`e il tasso ``bloccato oggi'' per il periodo
futuro $[t_1, t_2]$. Per assenza di arbitraggio, investire fino a $t_2$ deve dare lo stesso
risultato che investire fino a $t_1$ e poi reinvestire al tasso forward da $t_1$ a $t_2$:
\[
  P(0, t_2) = P(0, t_1) \cdot e^{-f(t_1,\,t_2)\,(t_2-t_1)},
\]
da cui, prendendo il logaritmo e isolando $f$:
\begin{equation}
  f(t_1, t_2) = \frac{\ln P(0,t_1) - \ln P(0,t_2)}{t_2 - t_1}.
  \label{eq:fwd-finito}
\end{equation}
Geometricamente, $f(t_1,t_2)$ \`e il \textbf{rapporto incrementale} (con segno) di
$\ln P(0,\cdot)$ sull'intervallo $[t_1,t_2]$: la pendenza media della curva
``log-sconto'' tra le due scadenze.

\paragraph{Forward istantaneo: limite per $t_2 \to t_1$.}

\begin{definizione}[Tasso forward istantaneo]
\label{def:forward}
Facendo tendere $t_2 \to t_1 = t$ nella~\eqref{eq:fwd-finito} si ottiene il
\textbf{tasso forward istantaneo}:
\begin{equation}
  f(t) = \lim_{t_2 \to t} f(t, t_2)
       = -\frac{d}{dt}\ln P(0,t)
       = -\frac{P'(0,t)}{P(0,t)},
  \label{eq:fwd-ist}
\end{equation}
cio\`e la \textbf{derivata del log-fattore di sconto} cambiata di segno.
\end{definizione}

$f(t)\,dt$ rappresenta il tasso applicato a un prestito infinitesimale stipulato in $t$
e rimborsato in $t+dt$: il forward istantaneo \`e il ``tasso marginale'' o ``tasso locale''
della struttura a termine.

% ------------------------------------------------------------
\subsection{La relazione che unifica tutto}
% ------------------------------------------------------------

Integrando la definizione~\eqref{eq:fwd-ist} si recupera il fattore di sconto:
\begin{equation}
  P(0,t) = \exp\!\left(-\int_0^t f(s)\,ds\right).
  \label{eq:B-da-fwd}
\end{equation}
Sostituendo~\eqref{eq:B-da-fwd} in~\eqref{eq:spot-inv} si ottiene la relazione
fondamentale tra spot e forward:
\begin{equation}
  r(t) = \frac{1}{t}\int_0^t f(s)\,ds.
  \label{eq:spot-media-fwd}
\end{equation}
\begin{keybox}[Lo spot \`e la media dei forward]
L'equazione~\eqref{eq:spot-media-fwd} dice che il \textbf{tasso spot $r(t)$ \`e la media
aritmetica del forward istantaneo} sull'intervallo $[0,t]$. Una media \emph{reagisce con
ritardo} rispetto al valore puntuale: se il forward cresce, lo spot cresce pi\`u lentamente;
se il forward converge all'$\UFR$ dall'alto (come in Figura~\ref{fig:motivazione}),
lo spot ci arriver\`a \emph{dopo}. Questo spiega perch\'e in figura il forward raggiunge
l'asintoto a scadenze sensibilmente pi\`u brevi dello spot.
\end{keybox}

\medskip
La tabella seguente riassume le tre rappresentazioni equivalenti.

\begin{center}
\renewcommand{\arraystretch}{1.4}
\begin{tabular}{lll}
\toprule
\textbf{Oggetto} & \textbf{Formula} & \textbf{Lettura} \\
\midrule
Fattore di sconto & $P(0,t)$ & valore oggi di $1$\,\texteuro{} in $t$ \\
Tasso spot (continuo) & $r(t) = -\ln P(0,t)/t$ & tasso medio da $0$ a $t$ \\
Tasso forward istant. & $f(t) = -\tfrac{d}{dt}\ln P(0,t)$ & tasso marginale in $t$ \\
Ricostruzione & $P(0,t) = e^{-\int_0^t f(s)\,ds}$ & $P$ da integrale del forward \\
\bottomrule
\end{tabular}
\end{center}

\subsection{Tassi par e dati di input}
\label{sec:inputdati}

Gli strumenti effettivamente quotati sul mercato non sono zero-coupon: i prezzi
direttamente osservabili sono quelli di contratti a cedola multipla. Il concetto chiave
\`e il \textbf{tasso par}.

Considera un'obbligazione con \textbf{valore nominale} $1$, cedola annua fissa $s$ pagata
alle date $j = 1, 2, \ldots, T$, e rimborso del nominale a scadenza $T$. Il suo
\textbf{prezzo} oggi \`e la somma dei valori attuali di ogni flusso:
\[
  \text{Prezzo}(s, T)
  = s \sum_{j=1}^{T} P(0,j) + P(0,T).
\]
Un titolo che quota \emph{alla pari} ha prezzo esattamente uguale al valore nominale.

\begin{definizione}[Tasso par]
\label{def:par}
Il \textbf{tasso par} $s_T$ di scadenza $T$ \`e il tasso cedola che rende il prezzo
dell'obbligazione uguale al suo valore nominale ($= 1$):
\[
  \text{Prezzo}(s_T, T) = 1
  \;\;\Longleftrightarrow\;\;
  s_T \sum_{j=1}^{T} P(0,j) + P(0,T) = 1
  \;\;\Longrightarrow\;\;
  s_T = \frac{1 - P(0,T)}{\displaystyle\sum_{j=1}^{T} P(0,j)}.
\]
\end{definizione}

\begin{osservazione}
Un \textbf{Interest Rate Swap} (IRS) a cedola annua con scadenza $T$ pu\`o essere visto
come lo scambio della gamba variabile (floating) con la gamba fissa (cedola $s$). Per
non-arbitrage la gamba variabile vale $1$ alla data di stipula (si ``rinnova'' a ogni data
cedola), quindi il tasso fisso che rende nullo il valore dello swap all'emissione \`e
\emph{esattamente} il tasso par della Definizione~\ref{def:par}. I contratti swap IRS sono
gli strumenti liquidi quotati sul mercato a cui applicare questa formula.
\end{osservazione}

I tassi swap di mercato incorporano tuttavia un piccolo premio per il rischio di
controparte, che deve essere rimosso prima della calibrazione (cfr.\ Sez.~\ref{sec:cra}).
Il dato di input per il metodo Smith--Wilson \`e quindi il \textbf{tasso par
after-CRA}:
\begin{equation}
  s_T^{\rm EIOPA} = \frac{1 - P(0,T)}{\sum_{j=1}^{T} P(0,j)} + \CRA .
  \label{eq:par}
\end{equation}
La relazione~\eqref{eq:par} lega i dati osservati (i par rate) all'oggetto da ricostruire
(la funzione di sconto $P(0,\cdot)$), ed \`e il \emph{vincolo} che la curva Smith--Wilson
dovr\`a soddisfare ai nodi di mercato.

\begin{keybox}[Dati di input: IRS EURIBOR]
Usiamo i par rate degli \textbf{Interest Rate Swap EUR contro EURIBOR~6M} (ticker Bloomberg
\texttt{EUSA<tenor> CMPL Curncy}), file \texttt{dati/tickersEUSA\_v3.xlsx}. Le scadenze sono
le \textbf{15 scadenze DLT ufficiali} EUR SWP valide dal 30~giugno al 31~dicembre~2025:
\[
  \{1,2,3,4,5,6,7,8,9,10,11,12,13,15,20\}\ \text{anni}.
\]
\end{keybox}

La Figura~\ref{fig:spot-fwd-par} mostra, per il mese di riferimento (dicembre~2025), lo spot
e il forward ricostruiti insieme ai 15 par rate di input.

\begin{figure}[H]
\centering
\includegraphics[width=0.92\textwidth]{../output/eiopa_rfr_sw_var/fig2_spot_forward_par.pdf}
\caption{Spot $r(0,T)$, forward $f(0,T)$ e par rate \texttt{EUSA*} (punti neri) ai 15 nodi DLT,
dicembre~2025. La curva Smith--Wilson interpola i par rate e produce uno spot e un forward
$\mathcal{C}^\infty$ su tutto $[0,\CP]$.}
\label{fig:spot-fwd-par}
\end{figure}

% ============================================================
\section{Perch\'e non bastano le spline: il problema dell'estrapolazione}
\label{sec:spline-limite}
% ============================================================

Disponiamo ora di tutto il necessario --- fattori di sconto, tassi spot e forward, par
rate di mercato (§\ref{sec:inputdati}) --- per la domanda operativa: come si costruisce la
curva continua che interpola i par rate osservati e la si prolunga oltre l'ultimo dato?
Lo strumento classico \`e la \textbf{spline cubica}: un polinomio a tratti di grado~3,
continuo fino alla derivata seconda, che passa esattamente per i nodi. \`E lo standard per
\emph{interpolare} dati racchiusi tra osservazioni note.

\paragraph{Dove la spline funziona e dove no.}
Il mercato EUR \`e liquido solo fino al Last Liquid Point ($\LLP = 20$ anni): oltre,
\textbf{non esistono dati}. Interpolare tra i nodi liquidi \`e legittimo; ma \emph{prolungare}
la curva per $T > \LLP$ non \`e pi\`u interpolazione, \`e \textbf{estrapolazione}. Una spline
non ha alcun ancoraggio economico fuori dall'intervallo dei dati: il suo andamento oltre
l'ultimo nodo dipende \emph{soltanto} dalla condizione al contorno scelta --- una convenzione
matematica --- non dai dati n\'e dall'economia.

La Figura~\ref{fig:spline-sw} lo rende evidente. Le due spline --- \textbf{naturale}
(derivata seconda nulla ai bordi) e \textbf{cubica standard} (\texttt{fmm}) --- usano gli
\emph{stessi} tassi liquidi ai 15 nodi DLT: nella zona liquida coincidono, perch\'e l\`i i
dati le determinano. Oltre il $\LLP$ divergono in modo \textbf{arbitrario} --- la cubica
\texttt{fmm} precipita addirittura verso tassi negativi --- al variare della sola
condizione al contorno: nessuna converge a un livello di equilibrio di lungo periodo.

\begin{figure}[H]
\centering
\includegraphics[width=0.92\textwidth]{../output/eiopa_rfr_sw_var/fig_spline_vs_sw.pdf}
\caption{Estrapolazioni a confronto, luglio~2025. Due spline cubiche (naturale e
\texttt{fmm}) costruite sugli stessi tassi liquidi ai 15 nodi DLT: coincidono per
$T \le \LLP$ ma divergono in modo arbitrario nella zona di estrapolazione (grigia),
nessuna convergendo all'$\UFR$. La curva Smith--Wilson (blu) \`e invece costruita attorno
all'$\UFR = 3.30\%$ (linea viola), verso cui converge.}
\label{fig:spline-sw}
\end{figure}

\begin{osservazione}[Serve un ancoraggio, non una convenzione]
La normativa di vigilanza non lascia il lungo termine alla scelta arbitraria di una
condizione al contorno: \textbf{impone} che il tasso forward converga all'$\UFR$ al
Convergence Point. Nessuna spline soddisfa questo requisito per costruzione. Occorre un
metodo che incorpori l'$\UFR$ \emph{dentro} la famiglia di curve ammissibili: \`e ci\`o che
fa Smith--Wilson (§\ref{sec:sw-metodo}), cercando la funzione di sconto come perturbazione
della curva asintotica $e^{-\omega t}$ e garantendo la convergenza all'$\UFR$ a priori.
\end{osservazione}

% ============================================================
\section{Struttura matematica del metodo Smith--Wilson}
\label{sec:sw-metodo}
% ============================================================

L'idea di Smith--Wilson \`e cercare la funzione di sconto in una \textbf{famiglia
parametrica} costruita attorno all'$\UFR$, scegliendo i parametri in modo da
\textbf{riprezzare esattamente} gli strumenti di mercato. La struttura \`e quella di un
problema di \emph{interpolazione con kernel}, che si formula come un problema di
\textbf{minima energia} (minima norma) risolubile con la sola algebra lineare.

Introduciamo la costruzione a partire dalla \emph{funzione di Wilson} $W$
(§\ref{sub:kernel-W}); riscrivendo $W(t,u)=e^{-\omega(t+u)}H(t,u)$ la riconduciamo subito alla
forma imperniata sul \emph{nucleo} $H$, su cui poggia tutta l'algebra del resto del capitolo.

%\subsection{Il kernel di Wilson --- formulazione con il nucleo \texorpdfstring{$H$}{H}}
%\label{sub:kernel-H}
%
%\begin{definizione}[Nucleo e funzione di Wilson, \cite{eiopa2025} sez.~9.7]
%Posto $\omega = \ln(1+\UFR)$ (intensit\`a dell'$\UFR$), il \textbf{nucleo (cuore) di Wilson}
%\`e la funzione simmetrica
%\[
%  H(t,u) = \alpha\,\min(t,u) - e^{-\alpha\,\max(t,u)}\,\sinh\!\bigl(\alpha\,\min(t,u)\bigr),
%  \qquad H(t,u) = H(u,t),
%\]
%dove $\alpha>0$ \`e un parametro di \emph{velocit\`a di convergenza}. Il nucleo $H$ \`e il
%\textbf{mattone} con cui si costruisce la curva: come vedremo (Sez.~\ref{sec:minenergia}), la
%sua matrice di Gram $[H(u_i,u_j)]$ \`e simmetrica definita positiva, propriet\`a su cui poggia
%l'intera costruzione algebrica \cite{fasshauer2007,lageras2016}.
%\end{definizione}
%
%Nella presentazione classica di EIOPA il nucleo compare moltiplicato per il fattore di sconto
%asintotico, dando la \textbf{funzione di Wilson}
%\[
%  W(t,u) = e^{-\omega(t+u)}\,H(t,u)
%         = e^{-\omega u}\,\underbrace{e^{-\omega t}H(t,u)}_{\displaystyle \psi_u(t)} .
%\]
%\`E la funzione $W$ che si visualizza naturalmente (Figura~\ref{fig:kernel}), perch\'e
%localizzata attorno al nodo $u$ e smorzata a zero per $t\to\infty$. In questa dispensa
%lavoriamo per\`o con il nucleo $H$: \`e esso a generare la matrice di Gram SPD su cui si fonda
%il metodo, e tenere separato il fattore $e^{-\omega t}$ rende trasparente l'algebra che segue.
%
%La funzione di sconto viene cercata nella forma
%\begin{equation}
%  P(t) = e^{-\omega t}\Bigl(1 + \textstyle\sum_{j} \zeta_j\, H(t,u_j)\Bigr)
%       = e^{-\omega t} + \textstyle\sum_{j} \zeta_j\,\psi_{u_j}(t),
%  \qquad \psi_{u_j}(t) := e^{-\omega t}H(t,u_j),
%  \label{eq:sw-P}
%\end{equation}
%cio\`e come \textbf{combinazione lineare} delle funzioni di base $\psi_{u_j}$ centrate nei nodi
%di pagamento $u_j$, pi\`u il termine di base $e^{-\omega t}$ che codifica la convergenza
%all'$\UFR$. Le $\psi_{u_j}$ coincidono, a meno del fattore costante $e^{-\omega u_j}>0$, con le
%funzioni di Wilson $W(\cdot,u_j)$ della Figura~\ref{fig:kernel}: stessa forma, stessa
%localizzazione. (Riscrivendo $P$ direttamente in termini di $W$ si otterrebbe l'espressione
%equivalente $P(t)=e^{-\omega t}+\sum_j \xi_j\,W(t,u_j)$ con $\xi_j = e^{\omega u_j}\zeta_j$:
%una semplice riparametrizzazione dei coefficienti.)
%
%I coefficienti $\zeta_j$ sono le incognite del problema: verranno determinati imponendo il
%riprezzamento degli strumenti di mercato (Sez.~\ref{sec:minenergia}). Si mostrer\`a che la
%soluzione di minima energia impone $\zeta = \Qmat b$, dove $b \in \mathbb{R}^n$ sono i
%\textbf{moltiplicatori di Lagrange} del vincolo di repricing, soluzione del
%sistema~\eqref{eq:sistema}. Nella formula finale si scrive quindi $(Qb)_j$ al posto di
%$\zeta_j$.
%
%\begin{figure}[H]
%\centering
%\includegraphics[width=0.88\textwidth]{../output/eiopa_rfr_sw_var/fig3_kernel_wilson.pdf}
%\caption{Funzioni di Wilson $W(t,u_j)=e^{-\omega(t+u_j)}H(t,u_j)$ centrate nei nodi
%$u_j \in \{1,5,10,20\}$. Ogni funzione \`e localizzata attorno al proprio nodo e si smorza a
%zero; a meno del fattore costante $e^{-\omega u_j}$ esse sono le funzioni di base $\psi_{u_j}$
%di cui la curva Smith--Wilson \`e combinazione lineare (eq.~\ref{eq:sw-P}).}
%\label{fig:kernel}
%\end{figure}

\subsection{La funzione di Wilson e la funzione di sconto}
\label{sub:kernel-W}

\begin{definizione}[Funzione di Wilson, \cite{eiopa2025} sez.~9.7]
Posto $\omega = \ln(1+\UFR)$ (intensit\`a dell'$\UFR$), la \textbf{funzione di Wilson} \`e
\[
  W(t,u) = e^{-\omega(t+u)}\Bigl[\alpha\,\min(t,u)
           - e^{-\alpha\,\max(t,u)}\,\sinh\!\bigl(\alpha\,\min(t,u)\bigr)\Bigr],
  \qquad W(t,u)=W(u,t),
\]
con $\alpha>0$ parametro di \emph{velocit\`a di convergenza}. \`E un \textbf{nucleo simmetrico
definito positivo}: la sua matrice di Gram $\mathbf{W}=[W(u_i,u_j)]$. Il fattore $e^{-\omega(t+u)}$ smorza $W$ a zero per
$t\to\infty$, da cui la convergenza della curva all'$\UFR$.
\end{definizione}

La funzione di sconto viene cercata nella forma
\begin{equation}
  P(t) = e^{-\omega t} + \textstyle\sum_{j} \xi_j\, W(t,u_j),
  \label{eq:sw-P-W}
\end{equation}
cio\`e come termine di base $e^{-\omega t}$ --- che codifica la convergenza all'$\UFR$ --- pi\`u
una \textbf{combinazione lineare} delle funzioni di Wilson $W(\cdot,u_j)$ centrate nei nodi di
pagamento $u_j$ (Figura~\ref{fig:kernel-W}). I coefficienti $\xi_j$ sono le \textbf{incognite}
del problema, da determinare imponendo che la curva ricostruita riprezzi gli strumenti
osservati sul mercato.

Ogni funzione di Wilson $W(\cdot,u_j)$ \`e localizzata attorno al proprio nodo $u_j$ e si smorza
a zero per $t\to\infty$: il suo peso diventa via via meno rilevante, cosicch\'e la curva
converge all'$\UFR$ (Figura~\ref{fig:kernel-W}).

\begin{figure}[H]
\centering
\includegraphics[width=0.88\textwidth]{../output/eiopa_rfr_sw_var/fig3_kernel_wilson.pdf}
\caption{Funzioni di Wilson $W(t,u_j)$ centrate nei nodi $u_j \in \{1,5,10,20\}$. Ogni funzione
\`e localizzata attorno al proprio nodo e si smorza a zero: la curva Smith--Wilson \`e la
combinazione lineare $\sum_j \xi_j\,W(\cdot,u_j)$ sommata al termine $e^{-\omega t}$
(eq.~\ref{eq:sw-P-W}).}
\label{fig:kernel-W}
\end{figure}

\paragraph{Dalla funzione di Wilson al nucleo \texorpdfstring{$H$}{H}.}
La funzione di Wilson contiene il fattore di sconto asintotico $e^{-\omega(t+u)}$ moltiplicato
per il \textbf{nucleo (cuore) di Wilson}, la funzione simmetrica
\[
  H(t,u) = \alpha\,\min(t,u) - e^{-\alpha\,\max(t,u)}\,\sinh\!\bigl(\alpha\,\min(t,u)\bigr),
  \qquad H(t,u) = H(u,t),
\]
cosicch\'e $W(t,u) = e^{-\omega(t+u)}H(t,u)$. Sostituendo questa identit\`a
nella~\eqref{eq:sw-P-W} e mettendo in evidenza il fattore comune $e^{-\omega t}$ si ottiene
\[
  P(t)
  = e^{-\omega t} + \sum_{j=1}^{m} \xi_j\, e^{-\omega(t+u_j)}\,H(t,u_j)
  = e^{-\omega t}\Bigl(1 + \sum_{j=1}^{m} \xi_j\, e^{-\omega u_j}\,H(t,u_j)\Bigr).
\]
Posto allora $\zeta_j := e^{-\omega u_j}\,\xi_j$ (una semplice \emph{riparametrizzazione} dei
coefficienti, con $e^{-\omega u_j}>0$), la funzione di sconto si riscrive nella forma compatta
imperniata sul nucleo $H$
\begin{equation}
  P(t) = e^{-\omega t}\Bigl(1 + \sum_{j=1}^{m} \zeta_j\, H(t,u_j)\Bigr),
  \label{eq:sw-P}
\end{equation}
dove $u_1,\dots,u_m$ sono gli $m=\LLP$ nodi di pagamento e $\zeta=(\zeta_1,\dots,\zeta_m)\transp$
\`e il nuovo vettore di coefficienti incogniti. \`E la forma~\eqref{eq:sw-P} che useremo nel
resto del capitolo: come vedremo, \`e il nucleo $H$ a generare la matrice di Gram simmetrica
definita positiva su cui poggia l'intera costruzione algebrica
\cite{fasshauer2007,lageras2016}.


\subsection{Il problema di minima energia: dalla curva al sistema lineare}
\label{sec:minenergia}

La funzione di sconto Smith--Wilson \`e individuata da un \textbf{problema di minimo
vincolato}: tra tutte le curve che riprezzano esattamente gli strumenti di mercato si
sceglie quella di \emph{energia minima}. Mostriamo come questo problema si imposti e si
risolva con i soli strumenti dell'algebra lineare --- matrici simmetriche definite positive
(SPD), forme quadratiche, prodotti scalari indotti da una matrice SPD, moltiplicatori di
Lagrange e proiezioni ortogonali --- e come conduca \textbf{esattamente} al sistema di
calibrazione. Procediamo individuando la \emph{funzione da minimizzare}, i \emph{vincoli} di mercato
e la \emph{soluzione} via moltiplicatori di Lagrange.

% ----------------------------------------------------------
\paragraph{La funzione da minimizzare: l'energia.}
% ----------------------------------------------------------
Fissato il kernel $H$, ricordiamo che $u_1,\ldots,u_m$ sono gli $m=\LLP=20$ nodi di
pagamento di EIOPA (anni 1 fino al Last Liquid Point) in cui la curva di sconto P(t) viene
determinata. La curva~\eqref{eq:sw-P} \`e individuata \emph{soltanto} dal
\textbf{vettore finito} di coefficienti $\zeta = (\zeta_1,\dots,\zeta_m)\transp \in
\mathbb{R}^m$: il problema \`e dunque finito-dimensionale, ambientato nel sottospazio
$\mathcal{V} = \operatorname{span}\{H(\cdot,u_j)\}_{j=1}^{m}$, e tutto ci\`o che segue \`e
algebra lineare su $\mathbb{R}^m$. Resta da scegliere \emph{quale} funzione di $\zeta$
minimizzare: il punto di partenza \`e un fatto generale di algebra lineare, che enunciamo con
precisione perch\'e \`e il cardine di tutta la costruzione.

\begin{osservazione}[Una matrice SPD induce un prodotto scalare]
\label{oss:spd-prodotto}
Sia $\mathbf{A}\in\mathbb{R}^{m\times m}$ \textbf{simmetrica definita positiva} (SPD).
Allora la forma $\langle x,y\rangle_{\mathbf{A}} := x\transp \mathbf{A}\, y$ \`e un
\textbf{prodotto scalare} su $\mathbb{R}^m$, perch\'e ne soddisfa i tre assiomi:
\begin{enumerate}
  \item \textbf{simmetria}: $\langle x,y\rangle_{\mathbf{A}} = \langle y,x\rangle_{\mathbf{A}}$,
        conseguenza di $\mathbf{A} = \mathbf{A}\transp$ (\,$\mathbf{A}$ \emph{simmetrica}\,);
  \item \textbf{bilinearit\`a}: $\langle \cdot,\cdot\rangle_{\mathbf{A}}$ \`e lineare in
        ciascun argomento, perch\'e tale \`e il prodotto matrice--vettore;
  \item \textbf{definita positivit\`a}: $\langle x,x\rangle_{\mathbf{A}} = x\transp\mathbf{A}\,x > 0$
        per ogni $x\neq 0$, conseguenza di $\mathbf{A}$ \emph{definita positiva}.
\end{enumerate}
Di conseguenza $\|x\|_{\mathbf{A}} := \sqrt{x\transp \mathbf{A}\, x}$ \`e una \textbf{norma}.
\end{osservazione}

\begin{osservazione}[La matrice di Wilson \`e SPD: l'energia]
\label{oss:energia}
La \textbf{matrice di Wilson} $\Hmat = [H(u_i,u_j)]_{i,j=1}^{m}$ \`e la matrice di Gram
dei nuclei $H(\cdot,u_j)$ centrati nei nodi di pagamento: poich\'e tali funzioni sono
linearmente indipendenti, $\Hmat$ \`e simmetrica e \textbf{definita positiva} (SPD).
Per l'Osservazione~\ref{oss:spd-prodotto}, $\Hmat$ induce su $\mathbb{R}^m$ il prodotto
scalare $\langle x,y\rangle_{\Hmat} = x\transp\Hmat\,y$ e la norma $\|x\|_{\Hmat} = \sqrt{x\transp\Hmat\,x}$.

Definiamo \textbf{energia} (o ruvidezza) della curva la forma quadratica
\[
  E(\zeta) = \zeta\transp \Hmat\, \zeta = \|\zeta\|_{\Hmat}^2 \ \ge 0 .
\]
\`E una misura di quanto i coefficienti $\zeta$ ``caricano'' i nuclei $H(\cdot,u_j)$,
e dipende \textbf{solo dalla matrice di Gram} $\Hmat$: minimizzarla seleziona la curva
pi\`u regolare tra quelle compatibili con i dati.
\end{osservazione}

% ----------------------------------------------------------
\paragraph{I vincoli di mercato (affini in $\zeta$).}
% ----------------------------------------------------------
Abbiamo $n$ IRS EUR contro EURIBOR~6M (ticker Bloomberg \texttt{EUSA*},
cfr.~Sez.~\ref{sec:inputdati}) alle scadenze DLT $T_1<\cdots<T_n\le\LLP$, ognuno quotato al
proprio par rate after-CRA $s_j$ (Def.~\ref{def:par}).

La condizione ``lo swap $j$ vale la pari'' (Def.~\ref{def:par}), esplicitando i flussi
annuali cedola + rimborso del nozionale, \`e:
\[
  s_j\,\sum_{i=1}^{T_j-1} P(u_i) + (1+s_j)\,P(u_{T_j}) = 1,
\]
cio\`e una \emph{somma pesata} dei fattori di sconto su tutti i $T_j$ nodi di pagamento dello
swap, non solo sulla scadenza. Con $n$ strumenti si hanno $n$ tali condizioni.

\begin{osservazione}[Cosa vale 1 e cosa no]
\label{oss:pari}
\`E cruciale non confondere due oggetti distinti.
\begin{itemize}
  \item I \textbf{fattori di sconto} $P(u_i)$ con $u_i > 0$ sono \emph{incogniti} e
        \emph{strettamente minori di 1}: $P(0)=1$, ma $P(u_i) < 1$ per $u_i>0$, perch\'e il
        denaro futuro vale meno di quello odierno.
  \item Il \textbf{prezzo dello swap} $j$ (scadenza $T_j$) --- la somma pesata dai flussi di
        cassa, $s_j\sum_{i=1}^{T_j}P(u_i)+P(u_{T_j})$ --- vale \emph{esattamente} 1
        perch\'e $s_j$ \`e scelto proprio cos\`i (Definizione~\ref{def:par}).
\end{itemize}
Conoscere i par rate $s_j$ (osservabili sul mercato) \textbf{non} ci d\`a direttamente i
$P(u_i)$ (incogniti). Le $n$ condizioni di tasso par sono $n$ equazioni lineari nelle $m$
incognite $P(u_1),\ldots,P(u_m)$; con $m = \LLP = 20$ nodi di pagamento e $n = 15$ scadenze
DLT, il sistema \`e \textbf{sottoindeterminato}: infiniti vettori $(P(u_i))$ soddisfano
simultaneamente le $n$ equazioni. Si nota inoltre che le scadenze $14, 16, 17, 18, 19$ anni
mancano dal panel DLT, rendendo inapplicabile anche un semplice bootstrap ricorsivo. La
matrice $\mathbf{C}$ scrive il sistema in forma compatta; il metodo Smith--Wilson seleziona
la soluzione di \emph{minima energia} (Sez.~\ref{sec:minenergia}).
\end{osservazione}

Per trattare simultaneamente le $n$ condizioni in forma algebrica compatta, raccogliamo i
flussi nella \textbf{matrice dei flussi di cassa} $\mathbf{C}\in\mathbb{R}^{m\times n}$,
dove la colonna $j$ \`e il profilo temporale dei pagamenti dello swap $j$ (scadenza $T_j$ anni):
\[
  C_{ij} =
  \begin{cases}
    s_j     & i < T_j \quad\text{(cedola),}\\
    1 + s_j & i = T_j \quad\text{(cedola + rimborso del nozionale),}\\
    0        & i > T_j \quad\text{(nessun flusso dopo la scadenza).}
  \end{cases}
\]
Per gli strumenti consecutivi $j=1,\ldots,13$ si ha $T_j=j$; per il 15Y e il 20Y invece
$T_{14}=15$ e $T_{15}=20$, per cui la riga del rimborso del nozionale non coincide con
l'indice di colonna.
La condizione di riprezzamento si scrive allora $\sum_{i=1}^{m} C_{ij}\,P(u_i) = 1$ per ogni
$j$ --- \`e la Definizione~\ref{def:par} riscritta esplicitando i flussi; in forma matriciale,
detto $p = \mathbf{1} \in \mathbb{R}^n$:
\[
  \mathbf{C}\transp \mathbf{P} = p .
\]

Valutiamo ora~\eqref{eq:sw-P} ai nodi per esprimere $\mathbf{P}$ in termini di $\zeta$.
Con $d_j = e^{-\omega u_j}$ si ha
\[
  \mathbf{P} = d + \operatorname{diag}(d)\,\Hmat\,\zeta .
\]
Sostituendo in $\mathbf{C}\transp \mathbf{P} = p$ e ponendo
$\Qmat = \operatorname{diag}(d)\,\mathbf{C}$ e $q = \mathbf{C}\transp d$, si ottiene la
\textbf{condizione di vincolo}
\begin{equation}
  \Qmat\transp \Hmat\, \zeta = p - q .
  \label{eq:vincolo}
\end{equation}
Si tratta di un \textbf{sistema lineare} di $n$ equazioni nelle $m$ incognite $\zeta$:
i suoi vettori soluzione formano un \textbf{insieme affine} (non vuoto) di $\mathbb{R}^m$.

% ----------------------------------------------------------
\paragraph{La soluzione: moltiplicatori di Lagrange.}
% ----------------------------------------------------------
Tra tutte le curve che riprezzano gli strumenti scegliamo quella di energia minima:
\begin{equation}
  \min_{\zeta\,\in\,\mathbb{R}^m} \ \zeta\transp \Hmat\, \zeta
  \qquad \text{s.t.}\quad
  \Qmat\transp \Hmat\, \zeta = p - q .
  \label{eq:minenergia}
\end{equation}
Questo problema \`e \textbf{ben posto} gi\`a prima di calcolarne la soluzione: la funzione
obiettivo $\zeta\mapsto\zeta\transp\Hmat\zeta$ \`e una forma quadratica
\textbf{strettamente convessa} (perch\'e $\Hmat$ \`e definita positiva) ristretta a un
insieme affine non vuoto. Un tale problema ammette uno e un solo punto di minimo.

\begin{proposizione}[Minima norma, versione finito-dimensionale]
\label{prop:minnorm-finita}
Il problema~\eqref{eq:minenergia} ammette \textbf{soluzione unica}. L'esistenza e
l'unicit\`a seguono dalla \textbf{definita positivit\`a} di $\Hmat$ (stretta
convessit\`a dell'energia) su un vincolo affine non vuoto.
\end{proposizione}

Per \emph{calcolare} il minimo si usano i \textbf{moltiplicatori di Lagrange}.
Ad ogni vincolo scalare $(\Qmat\transp\Hmat\zeta)_k = (p-q)_k$, $k=1,\ldots,n$,
si associa un moltiplicatore $b_k\in\mathbb{R}$; raccolti in $b\in\mathbb{R}^n$, si forma la
\textbf{funzione lagrangiana}
\begin{equation}
  \mathcal{L}(\zeta,b)
  = \zeta\transp\Hmat\zeta
    \;-\; 2\,b\transp\!\bigl(\Qmat\transp\Hmat\zeta - (p-q)\bigr).
  \label{eq:lagrangiana}
\end{equation}
Si calcola il gradiente rispetto a $\zeta$ e lo si pone uguale a zero:
\[
  \nabla_\zeta\,\mathcal{L}
  = 2\,\Hmat\zeta - 2\,\Hmat\Qmat\,b = 0
  \quad\Longrightarrow\quad
  \Hmat(\zeta - \Qmat b) = 0.
\]
Poich\'e $\Hmat$ \`e invertibile (SPD), si ottiene
\[
  \zeta = \Qmat\, b .
\]
Sostituendo nel vincolo~\eqref{eq:vincolo} si ritrova \textbf{esattamente} il sistema di
calibrazione~\eqref{eq:sistema}:
\[
  \bigl(\Qmat\transp \Hmat\, \Qmat\bigr)\, b = p - q .
\]
I coefficienti $b$ del metodo Smith--Wilson non sono dunque altro che i
\textbf{moltiplicatori di Lagrange} del problema di minima energia~\eqref{eq:minenergia}.

La relazione $\zeta = \Qmat\,b$ ha un significato \textbf{geometrico}: il vettore ottimo dei
coefficienti appartiene a $\operatorname{range}(\Qmat)$, il \textbf{sottospazio generato
dalle colonne di} $\Qmat$ (le sezioni del kernel pesate dai flussi di cassa). La soluzione
\`e quindi la \textbf{proiezione $\Hmat$-ortogonale} della curva base $e^{-\omega t}$
sull'insieme delle curve che riprezzano gli strumenti: tra tutte le correzioni ammissibili
si sceglie quella di lunghezza minima nella metrica indotta da $\Hmat$. \`E il principio di
\textbf{minima energia} nella sua forma pi\`u trasparente, e completamente elementare.

\begin{keybox}[Riepilogo: minima energia e proiezione]
La curva Smith--Wilson si ottiene con la sola algebra lineare:
\begin{enumerate}
  \item la matrice di Wilson $\Hmat$ \`e \textbf{SPD} $\Rightarrow$ definisce un prodotto
        scalare e una nozione di \emph{energia} $\zeta\transp\Hmat\zeta$;
  \item si minimizza l'energia sotto i \textbf{vincoli lineari} di repricing
        (\,\eqref{eq:minenergia}\,);
  \item i \textbf{moltiplicatori di Lagrange} $b$ risolvono il sistema SPD
        $(\Qmat\transp\Hmat\Qmat)b = p-q$, di ordine $n$ (Cholesky);
  \item geometricamente, la curva \`e la \textbf{proiezione ortogonale} (nella metrica
        $\Hmat$) sul sottospazio generato dalle funzioni $H$.
\end{enumerate}
Per approfondire: \cite{boyd2018} per i minimi quadrati vincolati e la soluzione di minima
norma via moltiplicatori di Lagrange, \cite{fasshauer2007} per l'interpolazione con kernel
definiti positivi, \cite{lageras2016} per il caso Smith--Wilson.
\end{keybox}

\subsection{Il sistema lineare di calibrazione}

Imporre il riprezzamento esatto degli strumenti conduce, con la matrice dei flussi di cassa
$\mathbf{C}$ introdotta nella formulazione dei vincoli (Sez.~\ref{sec:minenergia}) e con le
posizioni
$d_j = e^{-\omega u_j}$, $\Qmat = \mathrm{diag}(d)\,\mathbf{C}$,
$\Hmat = [H(u_i,u_j)]$, $q = \mathbf{C}\transp d$, al sistema lineare
(\cite{eiopa2025} sez.~9.8--9.10)
\begin{equation}
  \bigl(\Qmat\transp \Hmat\, \Qmat\bigr)\, b = p - q,
  \qquad p = \mathbf{1}\ \text{(par swap)} .
  \label{eq:sistema}
\end{equation}
Risolto per $b$, i coefficienti $Qb = \Qmat b$ entrano direttamente
nella~\eqref{eq:sw-P}. La matrice del sistema ha una struttura che ne garantisce la buona
risolubilit\`a numerica.

\begin{teorema}[Sistema SPD e fattorizzazione di Cholesky]
\label{thm:spd}
La matrice $\Qmat\transp \Hmat\, \Qmat \in \mathbb{R}^{n\times n}$ \`e \textbf{simmetrica
definita positiva}. Di conseguenza il sistema~\eqref{eq:sistema} ammette \textbf{soluzione
unica}, calcolabile in modo stabile tramite la \textbf{fattorizzazione di Cholesky}
$\Qmat\transp \Hmat\,\Qmat = R\transp R$ con $R$ triangolare superiore.
\end{teorema}

La Figura~\ref{fig:QHQ} visualizza la matrice $\Qmat\transp\Hmat\Qmat$ per il mese di
riferimento: simmetrica e con autovalore minimo strettamente positivo.

\begin{figure}[H]
\centering
\includegraphics[width=0.72\textwidth]{../output/eiopa_rfr_sw_var/fig4_matrice_QHQ.pdf}
\caption{Matrice del sistema $\Qmat\transp\Hmat\Qmat$ ($n=15$), dicembre~2025. \`E simmetrica
e definita positiva ($\lambda_{\min} > 0$): la fattorizzazione di Cholesky \`e applicabile ed
\`e circa due volte pi\`u efficiente di una LU generica.}
\label{fig:QHQ}
\end{figure}

\begin{teorema}[Rappresentazione Smith--Wilson]
\label{thm:rappresentazione}
Risolto il sistema~\eqref{eq:sistema}, la funzione di sconto~\eqref{eq:sw-P} riprezza
esattamente tutti gli strumenti di input ed \`e di classe $\mathcal{C}^\infty$ su $[0,\CP]$.
Da essa si ricavano spot e forward tramite
\[
  r(0,t) = P(t)^{-1/t}-1, \qquad
  f(0,t) = \omega - \frac{\sum_j (Qb)_j\,\partial_t H(t,u_j)}{1 + \sum_j (Qb)_j\, H(t,u_j)} .
\]
\end{teorema}

\subsection{Calibrazione di $\alpha$: un problema di ricerca di zeri}

Il parametro $\alpha$ governa la velocit\`a di convergenza all'$\UFR$. EIOPA lo fissa con un
criterio preciso (\cite{eiopa2025} sez.~9.14): si sceglie il \textbf{pi\`u piccolo}
$\alpha \ge \alpha_{\min}=0.05$ tale che il forward al Convergence Point disti dall'$\UFR$
meno di una tolleranza $\tau = 1$ bp:
\begin{equation}
  \alpha^\star = \min\Bigl\{\alpha \ge \alpha_{\min} : \bigl|f(0,\CP;\alpha) - \omega\bigr|
  \le \tau \Bigr\}.
  \label{eq:criterio-alpha}
\end{equation}
Definita $g(\alpha) = |f(0,\CP;\alpha) - \omega| - \tau$, la~\eqref{eq:criterio-alpha} \`e un
\textbf{problema di ricerca di zeri}, risolvibile con il metodo di \textbf{bisezione} (robusto)
o di \textbf{Newton} (rapido). La Figura~\ref{fig:alpha} mostra $g$ per il mese di riferimento.

\begin{figure}[H]
\centering
\includegraphics[width=0.9\textwidth]{../output/eiopa_rfr_sw_var/fig5_calibrazione_alpha.pdf}
\caption{Calibrazione di $\alpha$ come ricerca di zeri (dicembre~2025). La curva \`e
$f(0,\CP;\alpha)-\omega$ in bps; la banda grigia \`e la tolleranza $\pm 1$ bp. Il valore
$\alpha^\star \approx 0.064$ \`e il pi\`u piccolo $\alpha \ge \alpha_{\min}$ che entra nella
banda.}
\label{fig:alpha}
\end{figure}

\begin{teorema}[Convergenza all'UFR]
\label{thm:convergenza}
Per ogni $\alpha>0$ il forward Smith--Wilson converge all'intensit\`a dell'$\UFR$:
\[
  \lim_{t\to\infty} f(0,t) = \omega = \ln(1+\UFR).
\]
\end{teorema}

\begin{algobox}[Costruzione della curva Smith--Wilson (in R)]
\begin{Verbatim}[fontsize=\small]
# INPUT: par rate EUSA after-CRA r_k ai nodi T_mkt = {1..13,15,20}

# 1. Matrice dei flussi C (m x n), m = LLP date di pagamento annuali
C <- build_C(T_mkt, r)                 # C[j,k]=r_k (j<T_k), C[T_k,k]=1+r_k

# 2. Oggetti del sistema (dato alpha)
d   <- exp(-omega * (1:LLP)); Q <- C * d
H   <- outer(1:LLP, 1:LLP, H_heart, a = alpha)
QHQ <- t(Q) %*% H %*% Q                # SIMMETRICA DEFINITA POSITIVA

# 3. Soluzione via Cholesky (QHQ = R'R)
R  <- chol(QHQ)
b  <- backsolve(R, forwardsolve(t(R), p - q))
Qb <- as.numeric(Q %*% b)

# 4. alpha* via ricerca di zeri: |f(CP,alpha) - omega| <= 1bp, alpha >= 0.05
alpha_star <- alpha_star_for_mat(r, T_mkt)

# 5. Curva: P(t) = e^{-omega t}(1 + sum_j Qb_j H(t,u_j))
spot <- sw_spot_ann(v, Qb, alpha_star)
\end{Verbatim}
\end{algobox}

% ============================================================
\section{Ricostruzione passo a passo: luglio~2025}
\label{sec:ricostruzione}
% ============================================================

Percorriamo tutti i passi della teoria sul mese di \textbf{luglio~2025} (fine mese:
31 luglio~2025), usando i par rate \texttt{EUSA*} reali e confrontando il risultato
con la curva ufficiale pubblicata da EIOPA nel file \texttt{EIOPA\_RFR\_20250731.zip}.

% ----------------------------------------------------------
\subsection{Dati di input: par rate after-CRA}
% ----------------------------------------------------------

I par rate \texttt{EUSA*} (IRS EUR vs EURIBOR~6M) ai 15 nodi DLT, letti da
\texttt{dati/tickersEUSA\_v3.xlsx}, sono riportati nella Tabella~\ref{tab:input-jul}.
Il $\CRA = 10$ bps viene sottratto per ottenere i tassi after-CRA $r_j = s_j - \CRA$
che alimentano la calibrazione.

\begin{table}[H]
\centering\small
\caption{Par rate \texttt{EUSA*} e after-CRA per luglio~2025 (\%).}
\label{tab:input-jul}
\renewcommand{\arraystretch}{1.15}
\begin{tabular}{rcc}
\toprule
\textbf{Tenor} (anni) & $s_j$ = EUSA (\%) & $r_j = s_j - \CRA$ (\%)\\
\midrule
 1  & 2.0710 & 1.9710\\
 2  & 2.1008 & 2.0008\\
 3  & 2.1963 & 2.0963\\
 4  & 2.2840 & 2.1840\\
 5  & 2.3607 & 2.2607\\
 6  & 2.4340 & 2.3340\\
 7  & 2.4994 & 2.3994\\
 8  & 2.5628 & 2.4628\\
 9  & 2.6188 & 2.5188\\
10  & 2.6698 & 2.5698\\
11  & 2.7170 & 2.6170\\
12  & 2.7568 & 2.6568\\
13  & 2.7915 & 2.6915\\
15  & 2.8439 & 2.7439\\
20  & 2.8832 & 2.7832\\
\bottomrule
\end{tabular}
\end{table}

% ----------------------------------------------------------
\subsection{Costruzione degli oggetti del sistema}
% ----------------------------------------------------------

Con i tassi after-CRA $r_j$ e i nodi di pagamento $u_1=1,\dots,u_{20}=20$ si costruisce
la \textbf{matrice dei flussi di cassa} $\mathbf{C} \in \mathbb{R}^{20 \times 15}$
(Sez.~\ref{sec:minenergia}):
\[
  C_{ij} = \begin{cases}
    r_j     & i < T_j,\\
    1 + r_j & i = T_j,\\
    0       & i > T_j.
  \end{cases}
\]
Per esempio, la prima colonna (swap~1Y, $T_1=1$, $r_1=1.9710\%$):
\[
  \mathbf{c}_1 = (1 + 0.01971,\; 0,\; 0,\;\ldots,\; 0)^\top = (1.01971,\;0,\;\ldots)^\top.
\]
La colonna del 15Y (swap~15Y, $T_{14}=15$, $r_{14}=2.7439\%$) ha cedole
$r_{14} = 0.027439$ alle righe $i=1,\ldots,14$ e $1+r_{14}=1.027439$ alla riga $i=15$;
le righe $16,\ldots,20$ sono zero.

Da $\mathbf{C}$ si derivano gli oggetti del sistema (Sez.~\ref{sec:minenergia}):
\begin{itemize}
  \item $d_i = e^{-\omega u_i}$ con $\omega = \ln(1{+}\UFR)$;
  \item $\Qmat = \operatorname{diag}(d)\,\mathbf{C} \in \mathbb{R}^{20\times 15}$;
  \item $q_j = \sum_{i=1}^{20} d_i\, C_{ij}$ (colonna $j$ di $\Qmat$ sommata);
  \item $\Hmat = [H(u_i,u_j)]_{i,j=1}^{20}$ matrice SPD ($20\times 20$);
  \item $\Qmat^\top \Hmat\, \Qmat \in \mathbb{R}^{15\times 15}$ SPD (Teorema~\ref{thm:spd}).
\end{itemize}

% ----------------------------------------------------------
\subsection{Calibrazione di $\alpha$ e soluzione del sistema}
% ----------------------------------------------------------

Il parametro $\alpha$ \`e fissato con il criterio~\eqref{eq:criterio-alpha}: si cerca il
pi\`u piccolo $\alpha \ge 0.05$ tale che $|f(0,60;\alpha)-\omega| \le 1\,\text{bp}$.
La Figura~\ref{fig:alpha-jul} mostra la funzione $g(\alpha)$ per luglio~2025: lo zero
cade a $\alpha^\star = 0.0908$.

\begin{figure}[H]
\centering
\includegraphics[width=0.88\textwidth]{../output/eiopa_rfr_sw_var/fig_sw_alpha_jul.pdf}
\caption{Calibrazione di $\alpha$ con il criterio di ricerca di zeri, luglio~2025.
La curva \`e $g(\alpha)=f(0,\CP;\alpha)-\omega$ in bps; banda grigia = tolleranza
$\pm 1$ bp. Il valore $\alpha^\star = 0.0908$ \`e il pi\`u piccolo ammissibile.}
\label{fig:alpha-jul}
\end{figure}

Con $\alpha^\star = 0.0908$ si forma la matrice $\Qmat^\top\Hmat\Qmat$ e si risolve il
sistema (Sez.~\ref{sec:minenergia}) via fattorizzazione di Cholesky:
\[
  (\Qmat^\top \Hmat\, \Qmat)\, b = p - q,
  \qquad
  \zeta = \Qmat b.
\]

% ----------------------------------------------------------
\subsection{Curva calibrata e confronto con EIOPA}
% ----------------------------------------------------------

La Figura~\ref{fig:curve-jul} mostra spot e forward della curva Smith--Wilson calibrata
per luglio~2025. I punti neri sono i par rate \texttt{EUSA*}: la curva spot li interpola
esattamente (a meno di $\CRA$, già corretto). Il forward converge alla linea $\UFR$ per
$T > \LLP$.

\begin{figure}[H]
\centering
\includegraphics[width=0.92\textwidth]{../output/eiopa_rfr_sw_var/fig_sw_curve_jul.pdf}
\caption{Curva Smith--Wilson calibrata, luglio~2025. Spot (blu) e forward (rosso
mattone) fino a $T=30$ anni; punti neri = par rate EUSA* ai 15 nodi DLT; linea
viola = $\UFR = 3.30\%$.}
\label{fig:curve-jul}
\end{figure}

La Figura~\ref{fig:residui-jul} confronta la curva ricostruita con quella ufficiale
EIOPA su tutti i tenor interi da 1 a 80 anni, distinguendo la zona liquida ($T \le \LLP$)
dalla zona di estrapolazione ($T > \LLP$). In zona liquida gli scarti sono inferiori a
$1$ bp (RMSE $= 0.28$ bps); in estrapolazione la curva converge all'$\UFR$ e gli scarti
rimangono contenuti.

\begin{figure}[H]
\centering
\includegraphics[width=0.92\textwidth]{../output/eiopa_rfr_sw_var/fig_sw_residui_jul.pdf}
\caption{Scarti (SW $-$ ufficiale EIOPA) in bps su tutti i tenor interi 1--80 anni,
luglio~2025. Barre blu = zona liquida ($T \le \LLP$); barre arancio = estrapolazione.
Linee tratteggiate: $\pm 2$ bps.}
\label{fig:residui-jul}
\end{figure}

\begin{osservazione}
Il metodo Smith--Wilson riprezza esattamente gli strumenti di input (per costruzione),
ma la curva \emph{ufficiale} EIOPA utilizza dati da \textbf{Refinitiv}, mentre il nostro
script usa Bloomberg \texttt{EUSA*}: differenze di $\pm 0.1$ bp nei par rate si propagano
in scarti di pochi decimi di bp nella curva ricostruita. Il risultato conferma che
l'\textbf{IRS EURIBOR~6M} (\texttt{EUSA*}) \`e lo strumento di input coerente con la
curva ufficiale EIOPA.
\end{osservazione}

% ============================================================
\section{Conclusioni}
% ============================================================

\begin{itemize}
  \item La curva EIOPA risk-free risolve un problema di \textbf{interpolazione ed
        estrapolazione}: seguire i tassi di mercato fino al $\LLP$ e convergere all'$\UFR$ oltre.
  \item Il metodo \textbf{Smith--Wilson} costruisce la funzione di sconto come combinazione di
        \textbf{funzioni di base centrate nei nodi} (i nuclei di Wilson); la si caratterizza come curva di \textbf{minima energia}
        (minima norma) tra quelle che riprezzano i dati, in analogia con le spline.
  \item Il nucleo numerico \`e un \textbf{sistema lineare SPD} risolto con \textbf{Cholesky}
        (Teorema~\ref{thm:spd}) e una \textbf{ricerca di zeri} per il parametro $\alpha$
        (criterio sez.~9.14).
  \item La ricostruzione sui sette mesi giugno--dicembre~2025 riproduce la curva ufficiale EIOPA
        entro pochi bps, validando sia l'implementazione sia la scelta dell'input di mercato.
\end{itemize}

% ============================================================
\renewcommand{\refname}{Riferimenti bibliografici}
\begin{thebibliography}{9}

\bibitem{eiopa2025}
  \textbf{EIOPA} (2025).
  \textit{Technical documentation of the methodology to derive EIOPA's risk-free interest rate
  term structures}, EIOPA-BoS-25-599.
  Sezione~9 (Extrapolation and interpolation) per il metodo Smith--Wilson.

\bibitem{solvency2}
  \textbf{Direttiva 2009/138/CE} (Solvency~II) e
  \textbf{Regolamento Delegato (UE) 2015/35}.
  Quadro normativo per la valutazione market-consistent delle riserve tecniche.

\bibitem{smithwilson2001}
  \textbf{Smith, A., Wilson, T.} (2001).
  \textit{Fitting yield curves with long term constraints}.
  Bacon \& Woodrow Research Notes.

\bibitem{quarteroni2017}
  \textbf{Quarteroni, A., Saleri, F., Gervasio, P.} (2017).
  \textit{Calcolo Scientifico}, 6a ed. Springer.
  Fattorizzazione di Cholesky (cap.~5) e metodi per la ricerca di zeri (cap.~2).

\bibitem{burden2016}
  \textbf{Burden, R.L., Faires, J.D.} (2016).
  \textit{Numerical Analysis}, 10th ed. Cengage.
  Bisezione, Newton e sistemi lineari definiti positivi.

\bibitem{boyd2018}
  \textbf{Boyd, S., Vandenberghe, L.} (2018).
  \textit{Introduction to Applied Linear Algebra: Vectors, Matrices, and Least Squares}.
  Cambridge University Press.
  Cap.~16: minimi quadrati vincolati e soluzione di minima norma via moltiplicatori di
  Lagrange. Testo liberamente disponibile (Stanford/UCLA).

\bibitem{fasshauer2007}
  \textbf{Fasshauer, G.E.} (2007).
  \textit{Meshfree Approximation Methods with MATLAB}. World Scientific.
  Interpolazione con kernel definiti positivi: l'interpolante \`e combinazione lineare delle
  sezioni del kernel e i coefficienti risolvono il sistema con la matrice di Gram SPD.

\bibitem{lageras2016}
  \textbf{Lager\aa s, A., Lindholm, M.} (2016).
  Issues with the Smith--Wilson method.
  \textit{Insurance: Mathematics and Economics}, 68, 32--41.
  Trattazione del metodo Smith--Wilson come interpolazione con kernel.

\end{thebibliography}

\end{document}
